% ============================================================
% 3.1 CONTEXT
% Responsible member: Bao
% ============================================================

\section{Context}

\subsection{Background}

In recent years, film photography has experienced a significant revival worldwide, including in Vietnam. Increasing numbers of photography enthusiasts, artists, and collectors are returning to analog photography because of its distinctive aesthetic qualities, creative experience, and tangible value.

Unlike digital photography, film photography requires a more deliberate shooting process and produces unique visual characteristics that cannot always be replicated by digital filters.

As the analog photography community continues to grow, the demand for supporting services such as film developing, film scanning, photo printing, film preservation, and digital archiving has also increased. However, the service ecosystem surrounding film photography remains highly fragmented. Photographers often need to search for Film Labs through social media platforms, online communities, messaging applications, or personal recommendations.

Information regarding service availability, pricing, turnaround time, supported film formats, processing quality, and customer reviews is frequently scattered across multiple channels.

\subsection{Current Problems}

Despite the growing interest in film photography, several challenges remain:

\begin{itemize}
    \item Difficulty in discovering and comparing Film Labs.
    \item Lack of transparent information about pricing and services.
    \item Inconvenient service booking and order management.
    \item Limited real-time tracking of film processing orders.
    \item Fragmented communication between photographers and Film Labs.
    \item Difficulty in preserving and organizing digital film archives.
    \item Lack of a centralized community platform for sharing knowledge and experiences.
\end{itemize}

\subsection{Project Motivation}

The project is motivated by the need to create a centralized digital platform that connects film photography enthusiasts with Film Processing Labs.

By integrating service discovery, booking, order tracking, digital archiving, community interaction, and AI-powered assistance into a single platform, the system aims to simplify the overall film photography workflow.

The proposed platform will reduce the difficulty of finding reliable Film Labs while helping service providers manage their operations more efficiently. At the same time, photographers can access information, compare services, track orders, store digital scans, and interact with the film photography community through one unified system.

\subsection{Project Objective and Significance}

The main objective of this project is to develop an AI-powered platform that connects film photography enthusiasts with Film Processing Labs.

The platform aims to:

\begin{itemize}
    \item Connect photographers with suitable Film Labs.
    \item Provide transparent information about services and pricing.
    \item Support service booking and order management.
    \item Enable real-time order tracking.
    \item Provide digital film archive management.
    \item Support community interaction and knowledge sharing.
    \item Integrate AI-powered features to improve user experience.
\end{itemize}

The project is expected to contribute to the digital transformation of the film photography service ecosystem while preserving and supporting the growth of the analog photography community.